\ifdefined\gkver\else\def\gkver{student}\fi
\documentclass[\gkver]{gaokaozhenti}
\begin{document}

\chapter{2019年江苏理综}
%% paperType: 理综物理部分

\begin{enumerate}
\item
%% number: 1
%% paperName: 2019年江苏理综第 1 题 3 分
%% typeId: 1
%% type: 单选题
%% chapter: 交流电
%% point: 变压器
%% method: 无
%% score: 3
%% degree: 950
%% duplicateId: 0
%% body:
某理想变压器原、副线圈的匝数之比为 1:10，当输入电压增加 $20\UV$ 时，输出电压\xzanswer{D}
\fourchoices[answer=D]
{降低 $2\UV$}
{增加 $2\UV$}
{降低 $200\UV$}
{增加 $200\UV$}
%% answer: D
%% memo:
\memoanswer{
由理想变压器的规律 $\frac{\Delta U_{1}}{\Delta U_{2}} = \frac{n_{1}}{n_{2}}$ 可知，当 $\Delta U_{1} = 20\UV$ 时，有 $\frac{20\UV}{\Delta U_{2}} = \frac{1}{10}$，解得 $\Delta U_{2} = 200\UV$，故 D 正确。
}

\item
%% number: 2
%% paperName: 2019年江苏理综第 2 题 3 分
%% typeId: 1
%% type: 单选题
%% chapter: 力物体的平衡
%% point: 多力平衡
%% method: 无
%% score: 3
%% degree: 850
%% duplicateId: 0
%% body:
如图所示，一只气球在风中处于静止状态，风对气球的作用力水平向右。细绳与竖直方向的夹角为 $\alpha$，绳的拉力为 $T$，则风对气球作用力的大小为\xzanswer{C}
\onepicture[w=2.80cm]{figs/02.png}
\fourchoices[answer=C]
{$\frac{T}{\sin\alpha}$}
{$\frac{T}{\cos\alpha}$}
{$T\sin\alpha$}
{$T\cos\alpha$}
%% answer: C
%% memo:
\memoanswer{
气球在水平方向上仅受风对其的作用力与拉力的分力，由于气球静止水平方向受力平衡，由共点力平衡的规律知，风对气球作用力的大小为 $F = T\sin\alpha$，故 C 正确。
}

\item
%% number: 3
%% paperName: 2019年江苏理综第 3 题 3 分
%% typeId: 1
%% type: 单选题
%% chapter: 电路
%% point: 闭合电路欧姆定律
%% method: 无
%% score: 3
%% degree: 850
%% duplicateId: 0
%% body:
如图所示的电路中，电阻 $R = 2\UO$。断开 S 后，电压表的读数为 $3\UV$；闭合 S 后，电压表的读数为 $2\UV$，则电源的内阻 $r$ 为\xzanswer{A}
\onepicture[w=4.20cm]{figs/03.png}
\fourchoices[answer=A]
{$1\UO$}
{$2\UO$}
{$3\UO$}
{$4\UO$}
%% answer: A
%% memo:
\memoanswer{
断开 S 后，电源与电压表串联可知，电源电动势 $E = 3\UV$；闭合 S，$R$ 与 $r$ 串联，电压表测 $R$ 两端电压，由串联电路电流相等与闭合电路欧姆定律得 $\frac{3\UV - 2\UV}{r} = \frac{2\UV}{2\UO}$，解得 $r = 1\UO$，故 A 正确。
}

\item
%% number: 4
%% paperName: 2019年江苏理综第 4 题 3 分
%% typeId: 1
%% type: 单选题
%% chapter: 曲线运动
%% point: 人造地球卫星
%% method: 无
%% score: 3
%% degree: 650
%% duplicateId: 0
%% body:
1970 年成功发射的“东方红一号”是我国第一颗人造地球卫星，该卫星至今仍沿椭圆轨道绕地球运动。如图所示，设卫星在近地点、远地点的速度分别为 $v_{1}$、$v_{2}$，近地点到地心的距离为 $r$，地球质量为 $M$，引力常量为 $G$，则\xzanswer{B}
\onepicture[w=6.50cm]{figs/04.png}
\fourchoices[answer=B]
{$v_{1} > v_{2}$，$v_{1} = \sqrt{\frac{GM}{r}}$}
{$v_{1} > v_{2}$，$v_{1} > \sqrt{\frac{GM}{r}}$}
{$v_{1} < v_{2}$，$v_{1} = \sqrt{\frac{GM}{r}}$}
{$v_{1} < v_{2}$，$v_{1} > \sqrt{\frac{GM}{r}}$}
%% answer: B
%% memo:
\memoanswer{
如果卫星以半径 $r$ 绕地球做匀速圆周运动，由万有引力提供向心力有 $G\frac{Mm}{r^{2}} = m\frac{v^{2}}{r}$，则卫星的线速度 $v = \sqrt{\frac{GM}{r}}$；卫星在近地点的速度要大于以相同轨道半径做圆周运动的卫星的速度，才能到更高的椭圆轨道上运行，所以 $v_{1} > \sqrt{\frac{GM}{r}}$；卫星在轨运行时机械能守恒，其由近地点向远地点运行时，万有引力做负功，动能减小，所以 $v_{1} > v_{2}$，故 B 正确。
}

\item
%% number: 5
%% paperName: 2019年江苏理综第 5 题 3 分
%% typeId: 1
%% type: 单选题
%% chapter: 电场
%% point: 带电粒子的偏转
%% method: 无
%% score: 3
%% degree: 550
%% duplicateId: 0
%% body:
一匀强电场的方向竖直向上，$t = 0$ 时刻，一带电粒子以一定初速度水平射入该电场，电场力对粒子做功的功率为 $P$，不计粒子重力，则 $P-t$ 关系图象是\xzanswer{A}
\fourchoices[answer=A, ispicture=true, h=2.60cm]
{figs/05a.png}
{figs/05b.png}
{figs/05c.png}
{figs/05d.png}
%% answer: A
%% memo:
\memoanswer{
设竖直方向粒子速度大小为 $v$，加速度大小为 $a$，则电场力对粒子做功的功率为 $P = Fv = qEv = qEat = qE\frac{qE}{m}t = \frac{q^{2}E^{2}}{m}t$，则 $P$ 与 $t$ 成正比，故 A 正确。
}

\item
%% number: 6
%% paperName: 2019年江苏理综第 6 题 4 分
%% typeId: 2
%% type: 多选题
%% chapter: 曲线运动
%% point: 圆周运动的应用
%% method: 无
%% score: 4
%% degree: 750
%% duplicateId: 0
%% body:
如图所示，摩天轮悬挂的座舱在竖直平面内做匀速圆周运动。座舱的质量为 $m$，运动半径为 $R$，角速度大小为 $\omega$，重力加速度为 $g$，则座舱\xzanswer{BD}
\onepicture[w=4.20cm]{figs/06.png}
\fourchoices[answer=BD]
{运动周期为 $\frac{2\pi R}{\omega}$}
{线速度的大小为 $\omega R$}
{受摩天轮作用力的大小始终为 $mg$}
{所受合力的大小始终为 $m\omega^{2}R$}
%% answer: BD
%% memo:
\memoanswer{
根据匀速圆周运动的规律，座舱的运动周期为 $T = \frac{2\pi}{\omega}$，故 A 错误；座舱的线速度的大小为 $v = \omega R$，故 B 正确；座舱受摩天轮作用力和重力的合力等于向心力，故座舱受摩天轮作用力的大小始终不为 $mg$，故 C 错误；座舱做匀速圆周运动，所受合力的大小始终等于向心力，大小为 $m\omega^{2}R$，故 D 正确。
}

\item
%% number: 7
%% paperName: 2019年江苏理综第 7 题 4 分
%% typeId: 2
%% type: 多选题
%% chapter: 磁场
%% point: 左手定则
%% method: 无
%% score: 4
%% degree: 650
%% duplicateId: 0
%% body:
如图所示，在光滑的水平桌面上，a 和 b 是两条固定的平行长直导线，通过的电流强度相等。矩形线框位于两条导线的正中间，通有顺时针方向的电流，在 a、b 产生的磁场作用下静止。则 a、b 的电流方向可能是\xzanswer{CD}
\onepicture[w=4.20cm]{\includesvg{figs/07}}
\fourchoices[answer=CD]
{均向左}
{均向右}
{a 的向左，b 的向右}
{a 的向右，b 的向左}
%% answer: CD
%% memo:
\memoanswer{
导线 a、b 电流方向均向左，导线 a 对线框上面的边的安培力为斥力，导线 a 对线框下面的边的安培力为引力，导线 a 对线框的安培力的合力方向由 a 指向 b；导线 b 对线框上面的边的安培力为斥力，导线 b 对线框下面的边的安培力为引力，则导线 b 对线框的安培力的合力方向由 a 指向 b，故导线 a、b 对线框的合力不为零，故 A 错误；同理，导线 a、b 电流方向均向右，导线 a 对线框上面的边的安培力为引力，导线 a 对线框下面的边的安培力为斥力，导线 a 对线框的安培力的合力方向由 b 指向 a；导线 b 对线框上面的边的安培力为引力，导线 b 对线框下面的边的安培力为斥力，则导线 b 对线框的安培力的合力方向由 b 指向 a，故导线 a、b 对线框的合力不为零，故 B 错误；若导线 a 的电流方向向左，导线 b 的电流方向向右，导线 a 对线框上面的边的安培力为斥力，导线 a 对线框下面的边的安培力为引力，导线 a 对线框的安培力的合力方向由 a 指向 b，导线 b 对线框上面的边的安培力为引力，导线 b 对线框下面的边的安培力为斥力，导线 b 对线框的合力方向由 b 指向 a，且大小与导线 a 对线框的合力大小相等，故导线 a、b 对线框的合力为零，故 C 正确；同理，D 正确。
}

\item
%% number: 8
%% paperName: 2019年江苏理综第 8 题 4 分
%% typeId: 2
%% type: 多选题
%% chapter: 运动和力
%% point: 动量和能量
%% method: 无
%% score: 4
%% degree: 550
%% duplicateId: 0
%% body:
如图所示，轻质弹簧的左端固定，并处于自然状态。小物块的质量为 $m$，从 A 点向左沿水平地面运动，压缩弹簧后被弹回，运动到 A 点恰好静止。物块向左运动的最大距离为 $s$，与地面间的动摩擦因数为 $\mu$，重力加速度为 $g$，弹簧未超出弹性限度。在上述过程中\xzanswer{BC}
\onepicture[w=6.50cm]{figs/08.png}
\fourchoices[answer=BC]
{弹簧的最大弹力为 $\mu mg$}
{物块克服摩擦力做的功为 $2\mu mgs$}
{弹簧的最大弹性势能为 $\mu mgs$}
{物块在 A 点的初速度为 $\sqrt{2\mu gs}$}
%% answer: BC
%% memo:
\memoanswer{
小物块还能被弹回，说明最大弹力大于 $\mu mg$，故 A 错误；物块克服摩擦力做的功等于摩擦力乘以路程，即 $2\mu mgs$，故 B 正确；研究从弹簧被压缩到最短再到物块回到 A 点的过程，弹簧的最大弹性势能全部转化为内能，即 $\mu mgs$，故 C 正确；研究物块从 A 开始运动到弹簧压缩最短的过程中，由能量守恒得 $\frac{1}{2}mv_{A}^{2} = \mu mgs + E_{p}$，研究物块从 A 开始运动到最终停止的过程中，动能全部转化为内能，有 $\frac{1}{2}mv_{A}^{2} = 2\mu mgs$，解得 $v_{A} = 2\sqrt{\mu gs}$，故 D 错误。
}

\item
%% number: 9
%% paperName: 2019年江苏理综第 9 题 4 分
%% typeId: 2
%% type: 多选题
%% chapter: 电场
%% point: 电势能电势差电场力做功
%% method: 无
%% score: 4
%% degree: 450
%% duplicateId: 0
%% body:
如图所示，ABC 为等边三角形，电荷量为 $+q$ 的点电荷固定在 A 点。先将一电荷量也为 $+q$ 的点电荷 $Q_{1}$ 从无穷远处（电势为 0）移到 C 点，此过程中，电场力做功为 $-W$。再将 $Q_{1}$ 从 C 点沿 CB 移到 B 点并固定。最后将一电荷量为 $-2q$ 的点电荷 $Q_{2}$ 从无穷远处移到 C 点。下列说法正确的有\xzanswer{ABD}
\onepicture[w=4.20cm]{\includesvg{figs/09}}
\fourchoices[answer=ABD]
{$Q_{1}$ 移入之前，C 点的电势为 $\frac{W}{q}$}
{$Q_{1}$ 从 C 点移到 B 点的过程中，所受电场力做的功为 0}
{$Q_{2}$ 从无穷远处移到 C 点的过程中，所受电场力做的功为 $2W$}
{$Q_{2}$ 在移到 C 点后的电势能为 $-4W$}
%% answer: ABD
%% memo:
\memoanswer{
根据电势的定义，把 $Q_{1}$ 从 C 点移到无穷远处，电场力所做的功为 $W$，则 C 点的电势为 $\varphi = \frac{W}{q}$，故 A 正确；由于 B 点到 A 点的距离等于 C 点到 A 点的距离，根据点电荷的电势分布可知，B、C 点电势相等，$Q_{1}$ 从 C 点移到 B 点的过程中，所受电场力做的功为 $W_{1} = q(\varphi_{C} - \varphi_{B}) = 0$，故 B 正确；$Q_{2}$ 从无穷远处移到 C 点的过程中，由对称性可知，$Q_{1}$ 固定在 B 点后，C 点的电势为 $\frac{2W}{q}$，则 $Q_{2}$ 在 C 点的电势能 $E_{p} = -2q\varphi_{C} = -2q \cdot \frac{2W}{q} = -4W$，即 $Q_{2}$ 从无穷远处移到 C 点的过程中，电场力做功 $4W$，故 C 错误、D 正确。
}

\item
%% number: 10
%% paperName: 2019年江苏理综第 10 题 8 分
%% typeId: 4
%% type: 实验题
%% chapter: 功和能
%% point: 验证动能定理
%% method: 无
%% score: 8
%% degree: 550
%% duplicateId: 0
%% body:
某兴趣小组用如题 1 图所示的装置验证动能定理。
\onepicture[w=8.00cm, align=right]{figs/10a.png}
\begin{enumerate}
	\item
	有两种工作频率均为 $50\UHz$ 的打点计时器供实验选用：A．电磁打点计时器；B．电火花打点计时器。为使纸带在运动时受到的阻力较小，应选择\tkanswer{}（选填“A”或“B”）。
	\item
	保持长木板水平，将纸带固定在小车后端，纸带穿过打点计时器的限位孔。实验中，为消除摩擦力的影响，在砝码盘中慢慢加入沙子，直到小车开始运动。同学甲认为此时摩擦力的影响已得到消除。同学乙认为还应从盘中取出适量沙子，直至轻推小车观察到小车做匀速运动。看法正确的同学是\tkanswer{}（选填“甲”或“乙”）。
	\item
	消除摩擦力的影响后，在砝码盘中加入砝码。接通打点计时器电源，松开小车，小车运动。纸带被打出一系列点，其中的一段如题 2 图所示。图中纸带按实际尺寸画出，纸带上 A 点的速度 $v_{A} =$ \tkanswer{} $\Ums$。
	
	\onepicture[w=7.00cm, align=right]{figs/10b.png}
	\item
	测出小车的质量为 $M$，再测出纸带上起点到 A 点的距离为 $L$。小车动能的变化量可用 $\Delta E_{k} = \frac{1}{2}Mv_{A}^{2}$ 算出。砝码盘中砝码的质量为 $m$，重力加速度为 $g$；实验中，小车的质量应\tkanswer{}（选填“远大于”“远小于”或“接近”）砝码、砝码盘和沙子的总质量，小车所受合力做的功可用 $W = mgL$ 算出。多次测量，若 $W$ 与 $\Delta E_{k}$ 均基本相等则验证了动能定理。
\end{enumerate}
%% answer:
\jdanswer{
\begin{enumerate}
	\item
	B
	
	\item
	乙
	
	\item
	0.31（$0.30\sim0.33$ 都算对）
	
	\item
	远大于
\end{enumerate}
}
%% memo:
\memoanswer{
（1）电火花打点计时器的阻力较小，故 B 正确。

（2）甲同学做法抵消的是最大静摩擦力，而实际情况下滑动摩擦力比最大静摩擦力略小，乙同学的看法正确。

（3）在 A 点的左右两边取关于 A 对称的点（左、右各 4 个点以内），测出两点间的距离 $\Delta x$，由打点计时器可以得到所选两点间的时间间隔 $\Delta t$，则可以求出 A 点的速度 $v = \frac{\Delta x}{\Delta t} = 0.31\Ums$。

（4）实验中，用砝码的重力当作小车的合力，所以砝码的质量应该远小于小车的质量。
}

\item
%% number: 11
%% paperName: 2019年江苏理综第 11 题 10 分
%% typeId: 4
%% type: 实验题
%% chapter: 电路
%% point: 电阻定律
%% method: 无
%% score: 10
%% degree: 550
%% duplicateId: 0
%% body:
某同学测量一段长度已知的电阻丝的电阻率。实验操作如下：
\begin{enumerate}
	\item
	螺旋测微器如题 1 图所示。在测量电阻丝直径时，先将电阻丝轻轻地夹在测砧与测微螺杆之间，再旋动\tkanswer{}（选填“A”“B”或“C”），直到听见“喀喀”的声音，以保证压力适当，同时防止螺旋测微器的损坏。
	
	\onepicture[w=6.00cm, align=right]{figs/11a.png}
	\item
	选择电阻丝的\tkanswer{}（选填“同一”或“不同”）位置进行多次测量，取其平均值作为电阻丝的直径。
	\item
	2 图甲中 $R_{x}$ 为待测电阻丝。请用笔画线代替导线，将滑动变阻器接入 2 图乙实物电路中的正确位置\tkanswer{}。
	
	\twopicture[numstyle=chinese, wA=3.60cm, wB=6.40cm, align=right]{figs/11b.png}{figs/11c.png}
	\item
	为测量 $R$，利用 2 图甲所示的电路，调节滑动变阻器测得 5 组电压 $U_{1}$ 和电流 $I_{1}$ 的值，作出 $U_{1}-I_{1}$ 关系图象如图所示。接着，将电压表改接在 a、b 两端，测得 5 组电压 $U_{2}$ 和电流 $I_{2}$ 的值，数据见下表：
	
	\begin{center}
	\begin{tabular}{|c|c|c|c|c|c|}
	\hline
	$U_{2}/\UV$ & 0.50 & 1.02 & 1.54 & 2.05 & 2.55 \\
	\hline
	$I_{2}/\UmA$ & 20.0 & 40.0 & 60.0 & 80.0 & 100.0 \\
	\hline
	\end{tabular}
	\end{center}
	
	请根据表中的数据，在方格纸上作出 $U_{2}-I_{2}$ 图象。
	
	\onepicture[w=6.50cm, align=right]{figs/11d.png}
	\item
	由此，可求得电阻丝的 $R_{x} =$ \tkanswer{} $\UO$，根据电阻定律可得到电阻丝的电阻率。
\end{enumerate}
%% answer:
\jdanswer{
\begin{enumerate}
	\item
	C
	
	\item
	不同
	
	\item
	实物连线如答图甲所示。
	
	\item
	图象如答图乙所示。
	
	\item
	23.5（$23.0\sim24.0$ 都算对）
\end{enumerate}
}
%% memo:
\memoanswer{
（1）先将电阻丝轻轻地夹在测砧与测微螺杆之间，再旋动微调螺丝 C。

（2）为了减小误差，选择电阻丝的不同位置进行多次测量，取其平均值作为电阻丝的直径。

（3）根据电路图，连接实物图，如答图甲所示。

\onepicture[num=甲, labelprefix={答图}, w=5.00cm]{figs/11_an1.png}

（4）根据表格中的数据描点连线，如答图乙所示。

\onepicture[num=乙, labelprefix={答图}, w=6.00cm]{figs/11_an2.png}

（5）结合所画出的图象，可求得电阻丝的阻值 $R_{x} = \frac{\Delta U_{1}}{\Delta I_{1}} - \frac{\Delta U_{2}}{\Delta I_{2}} = 23.5\UO$。
}

\item
%% number: 12
%% paperName: 2019年江苏理综第 12 题 4 分
%% typeId: 6
%% type: 计算题
%% chapter: 光学
%% point: 光子说
%% method: 无
%% score: 4
%% degree: 550
%% duplicateId: 0
%% body:
本题为选修 3-5 模块。
\begin{enumerate}
	\item
	质量为 $M$ 的小孩站在质量为 $m$ 的滑板上，小孩和滑板均处于静止状态，忽略滑板与地面间的摩擦。小孩沿水平方向跃离滑板，离开滑板时的速度大小为 $v$，此时滑板的速度大小为\xzanswer{B}
	
	\fourchoices[answer=B]
	{$\frac{m}{M}v$}
	{$\frac{M}{m}v$}
	{$\frac{m}{m+M}v$}
	{$\frac{M}{m+M}v$}
	\item
	100 年前，卢瑟福用 $\alpha$ 粒子轰击氮核打出了质子。后来，人们用 $\alpha$ 粒子轰击 \ce{^{60}_{28}Ni} 核也打出了质子：\ce{^{4}_{2}He + ^{60}_{28}Ni -> ^{62}_{29}Cu + ^{1}_{1}H + X}；该反应中的 X 是\tkanswer{}（选填“电子”“正电子”或“中子”）。此后，对原子核反应的持续研究为核能利用提供了可能。目前人类获得核能的主要方式是\tkanswer{}（选填“核衰变”“核裂变”或“核聚变”）。
	\item
	在“焊接”视网膜的眼科手术中，所用激光的波长 $\lambda = 6.4\times10^{-7}\Um$，每个激光脉冲的能量 $E = 1.5\times10^{-2}\UJ$。求每个脉冲中的光子数目。（已知普朗克常量 $h = 6.63\times10^{-34}\UJcdots$，光速 $c = 3\times10^{8}\Ums$，计算结果保留一位有效数字）
\end{enumerate}
%% answer:
\jdanswer{
\begin{enumerate}
	\item
	B
	
	\item
	中子，核裂变
	
	\item
	$5\times10^{16}$ 个
\end{enumerate}
}
%% memo:
\memoanswer{
（1）小孩和滑板组成的系统在水平方向上不受外力，故水平方向系统动量守恒，令小孩离开滑板后滑板的速度为 $v'$，据动量守恒定律有 $Mv - mv' = 0$，解得 $v' = \frac{Mv}{m}$，故 B 正确。

（2）由核反应过程电荷数和质量数守恒可得，X 的质量数为 $60 + 4 - 62 - 1 = 1$，电荷数为 $2 + 28 - 29 - 1 = 0$，可以判断 X 是中子 \ce{^{1}_{0}n}；目前人类获得核能的主要方式是核裂变。

（3）光子能量 $\varepsilon = h\nu = \frac{hc}{\lambda}$，光子数目 $n = \frac{E}{\varepsilon}$，代入数据得 $n = 5\times10^{16}$ 个。
}

\item
%% number: 13
%% paperName: 2019年江苏理综第 13 题 4 分
%% typeId: 6
%% type: 计算题
%% chapter: 光学
%% point: 全反射
%% method: 无
%% score: 4
%% degree: 650
%% duplicateId: 0
%% body:
本题包括 A、B 两小题，请选定其中一小题，并在相应的答题区域内作答。若多做，则按 A 小题评分。

\textbf{A．}（选修模块 3-3）

\begin{enumerate}
	\item
	在没有外界影响的情况下，密闭容器内的理想气体静置足够长时间后，该气体\xzanswer{CD}
	
	\fourchoices[answer=CD]
	{分子的无规则运动停息下来}
	{每个分子的速度大小均相等}
	{分子的平均动能保持不变}
	{分子的密集程度保持不变}
	\item
	由于水的表面张力，荷叶上的小水滴总是球形的。在小水滴表面层中，水分子之间的相互作用总体上表现为\tkanswer{}（选填“引力”或“斥力”）。分子势能 $E_{p}$ 和分子间距离 $r$ 的关系图象如图所示，能总体上反映小水滴表面层中水分子 $E_{p}$ 的是图中\tkanswer{}（选填“A”“B”或“C”）的位置。
	
	\onepicture[w=4.80cm, align=right]{figs/13a.png}
	\item
	如图所示，一定质量理想气体经历 $A\rightarrow B$ 的等压过程，$B\rightarrow C$ 的绝热过程（气体与外界无热量交换），其中 $B\rightarrow C$ 过程中内能减少 $900\UJ$。求 $A\rightarrow B\rightarrow C$ 过程中气体对外界做的总功。
	
	\onepicture[w=5.00cm, align=right]{figs/13b.png}
\end{enumerate}

\textbf{B．}（选修模块 3-4）

\begin{enumerate}
	\item
	一单摆做简谐运动，在偏角增大的过程中，摆球的\xzanswer{AC}
	
	\fourchoices[answer=AC]
	{位移增大}
	{速度增大}
	{回复力增大}
	{机械能增大}
	\item
	将两支铅笔并排放在一起，中间留一条狭缝，通过这条狭缝去看与其平行的日光灯，能观察到彩色条纹，这是由于光的\tkanswer{}（选填“折射”“干涉”或“衍射”）。当缝的宽度\tkanswer{}（选填“远大于”或“接近”）光波的波长时，这种现象十分明显。
	\item
	如图所示，某 L 形透明材料的折射率 $n = 2$。现沿 AB 方向切去一角，AB 与水平方向的夹角为 $\theta$。为使水平方向的光线射到 AB 面时不会射入空气，求 $\theta$ 的最大值。
	
	\onepicture[w=5.60cm, align=right]{figs/13c.png}
\end{enumerate}
%% answer:
\jdanswer{
\begin{enumerate}
	\item
	A：（1）CD；（2）引力，C；（3）$W = 1500\UJ$。
	
	\item
	B：（1）AC；（2）衍射，接近；（3）$\theta = 60^{\circ}$。
\end{enumerate}
}
%% memo:
\memoanswer{
\textbf{A．}（1）分子永不停息地做无规则运动，不会停息下来，且每个分子运动速度大小不一定相等，选项 AB 错误；稳定之后，温度不变，分子的平均动能保持不变，选项 C 正确；理想气体质量不变，体积不变，所以密度不变，即分子的密集程度保持不变，选项 D 正确。

（2）在小水滴表面层中，相邻水分子之间的距离大于平衡距离 $r_{0}$，相互作用总体上表现为引力；图中 $E_{p}$ 最小的位置 B 点是 $r_{0}$ 处，所以能总体上反映小水滴表面层中水分子 $E_{p}$ 的是题图 1 中 C 的位置。

（3）$A\rightarrow B$ 过程气体所做的功为 $W_{1} = -p(V_{B} - V_{A})$，$B\rightarrow C$ 过程，由热力学第一定律得 $W_{2} = \Delta U$，则气体对外界做的总功为 $W = -(W_{1} + W_{2})$，代入数据得 $W = 1500\UJ$。

\textbf{B．}（1）在偏角增大的过程中，摆球的位移增大，摆球速度减小，回复力 $F = -kx$ 增大，小球的机械能守恒，故 AC 正确，BD 错误。

（2）通过一条狭缝观察到彩色条纹，这是光的衍射现象；缝的宽度越接近光波的波长，衍射现象越明显。

（3）水平方向光线刚好不从 AB 射出时，$\theta$ 角最大，即入射光刚好发生全反射，有 $\sin C = \frac{1}{n}$，且 $C + \theta = 90^{\circ}$，代入数据解得 $\theta = 60^{\circ}$。
}

\item
%% number: 14
%% paperName: 2019年江苏理综第 14 题 15 分
%% typeId: 6
%% type: 计算题
%% chapter: 电磁感应
%% point: 电磁感应的电量
%% method: 无
%% score: 15
%% degree: 650
%% duplicateId: 0
%% body:
如图所示，匀强磁场中有一个用软导线制成的单匝闭合线圈，线圈平面与磁场垂直。已知线圈的面积 $S = 0.3\Umq$、电阻 $R = 0.6\UO$，磁场的磁感应强度 $B = 0.2\UT$。现同时向两侧拉动线圈，线圈的两边在 $\Delta t = 0.5\Us$ 时间内合到一起。求线圈在上述过程中：
\begin{enumerate}
	\item
	感应电动势的平均值 $E$；
	\item
	感应电流的平均值 $I$，并在图中标出电流方向；
	\item
	通过导线横截面的电荷量 $q$。
\end{enumerate}
\onepicture[w=5.50cm, align=right]{figs/14.png}
%% answer:
\jdanswer{
\begin{enumerate}
	\item
	$E = 0.12\UV$
	
	\item
	$I = 0.2\UA$（电流方向见图）
	
	\item
	$q = 0.1\UC$
\end{enumerate}
}
%% memo:
\memoanswer{
（1）由法拉第电磁感应定律得 $E = \frac{\Delta\Phi}{\Delta t}$，磁通量的变化量为 $\Delta\Phi = B\Delta S$，解得 $E = \frac{B\Delta S}{\Delta t}$，代入数据得 $E = 0.12\UV$。

（2）回路中的电流大小为 $I = \frac{E}{R}$，代入数据得 $I = 0.2\UA$，方向如图所示。

\onepicture[w=4.00cm]{figs/14_an.png}

（3）由电荷量的表达式得 $q = I\Delta t$，解得 $q = 0.1\UC$。
}

\item
%% number: 15
%% paperName: 2019年江苏理综第 15 题 16 分
%% typeId: 6
%% type: 计算题
%% chapter: 运动和力
%% point: 牛顿定律的应用
%% method: 无
%% score: 16
%% degree: 450
%% duplicateId: 0
%% body:
如图所示，质量相等的物块 A 和 B 叠放在水平地面上，左边缘对齐。A 与 B、B 与地面间的动摩擦因数均为 $\mu$。先敲击 A，A 立即获得水平向右的初速度，在 B 上滑动距离 $L$ 后停下。接着敲击 B，B 立即获得水平向右的初速度，A、B 都向右运动，左边缘再次对齐时恰好相对静止，此后两者一起运动至停下。最大静摩擦力等于滑动摩擦力，重力加速度为 $g$。求：
\begin{enumerate}
	\item
	A 被敲击后获得的初速度大小 $v_{A}$；
	\item
	在左边缘再次对齐的前、后，B 运动加速度的大小 $a_{B}$、$a_{B}'$；
	\item
	B 被敲击后获得的初速度大小 $v_{B}$。
\end{enumerate}
\onepicture[w=6.00cm, align=right]{\includesvg{figs/15}}
%% answer:
\jdanswer{
\begin{enumerate}
	\item
	$v_{A} = \sqrt{2\mu gL}$
	
	\item
	$a_{B} = 3\mu g$，$a_{B}' = \mu g$
	
	\item
	$v_{B} = 2\sqrt{2\mu gL}$
\end{enumerate}
}
%% memo:
\memoanswer{
（1）A 被敲击后，B 静止。由牛顿运动定律知，A 加速度的大小 $a_{A} = \mu g$。A 做匀变速直线运动，由运动学公式得 $2a_{A}L = v_{A}^{2}$，解得 $v_{A} = \sqrt{2\mu gL}$。

（2）设 A、B 的质量均为 $m$。对齐前，B 所受合外力大小 $F = 3\mu mg$，由牛顿运动定律 $F = ma_{B}$，得 $a_{B} = 3\mu g$；对齐后，A、B 所受合外力大小 $F' = 2\mu mg$，由牛顿运动定律 $F' = 2ma_{B}'$，得 $a_{B}' = \mu g$。

（3）经过时间 $t$，A、B 达到共同速度 $v$，位移分别为 $x_{A}$、$x_{B}$，A 加速度的大小等于 $a_{A}$，则 $v = a_{A}t$，$v = v_{B} - a_{B}t$，$x_{A} = \frac{1}{2}a_{A}t^{2}$，$x_{B} = v_{B}t - \frac{1}{2}a_{B}t^{2}$，且 $x_{B} - x_{A} = L$，解得 $v_{B} = 2\sqrt{2\mu gL}$。
}

\item
%% number: 16
%% paperName: 2019年江苏理综第 16 题 16 分
%% typeId: 6
%% type: 计算题
%% chapter: 磁场
%% point: 带电粒子在磁场中的运动
%% method: 无
%% score: 16
%% degree: 350
%% duplicateId: 0
%% body:
如图所示，匀强磁场的磁感应强度大小为 $B$。磁场中的水平绝缘薄板与磁场的左、右边界分别垂直相交于 M、N，MN = $L$，粒子打到板上时会被反弹（碰撞时间极短），反弹前后水平分速度不变，竖直分速度大小不变、方向相反。质量为 $m$、电荷量为 $-q$ 的粒子速度一定，可以从左边界的不同位置水平射入磁场，在磁场中做圆周运动的半径为 $d$，且 $d < L$，粒子重力不计，电荷量保持不变。
\begin{enumerate}
	\item
	求粒子运动速度的大小 $v$；
	\item
	欲使粒子从磁场右边界射出，求入射点到 M 的最大距离 $d_{m}$；
	\item
	从 P 点射入的粒子最终从 Q 点射出磁场，PM = $d$，QN = $\frac{d}{2}$，求粒子从 P 到 Q 的运动时间 $t$。
\end{enumerate}
\onepicture[w=7.00cm, align=right]{figs/16.png}
%% answer:
\jdanswer{
\begin{enumerate}
	\item
	$v = \frac{qBd}{m}$
	
	\item
	$d_{m} = \frac{2+\sqrt{3}}{2}d$
	
	\item
	当 $L = nd + \left(1 - \frac{\sqrt{3}}{2}\right)d$ 时，$t = \left(\frac{L}{d} + \frac{3\sqrt{3}-4}{6}\right)\frac{\pi m}{2qB}$；当 $L = nd + \left(1 + \frac{\sqrt{3}}{2}\right)d$ 时，$t = \left(\frac{L}{d} - \frac{3\sqrt{3}-4}{6}\right)\frac{\pi m}{2qB}$。
\end{enumerate}
}
%% memo:
\memoanswer{
（1）粒子在磁场中做匀速圆周运动，由洛伦兹力提供向心力有 $qvB = m\frac{v^{2}}{R}$，解得 $R = \frac{mv}{qB}$，粒子的运动半径为 $d$，解得 $v = \frac{qBd}{m}$。

（2）如图所示，粒子碰撞后的运动轨迹恰好与磁场左边界相切，由几何关系得 $d_{m} = d(1 + \sin 60^{\circ})$，解得 $d_{m} = \frac{2+\sqrt{3}}{2}d$。

\onepicture[w=4.00cm]{figs/16_an.png}

（3）粒子的运动周期 $T = \frac{2\pi m}{qB}$。设粒子最后一次碰撞到射出磁场的时间为 $t'$，则 $t = n\frac{T}{4} + t'$（$n = 1,3,5,\cdots$）。当 $L = nd + \left(1 - \frac{\sqrt{3}}{2}\right)d$ 时，粒子斜向上射出磁场，$t' = \frac{1}{12}T$，解得 $t = \left(\frac{L}{d} + \frac{3\sqrt{3}-4}{6}\right)\frac{\pi m}{2qB}$；当 $L = nd + \left(1 + \frac{\sqrt{3}}{2}\right)d$ 时，粒子斜向下射出磁场，$t' = \frac{5}{12}T$，解得 $t = \left(\frac{L}{d} - \frac{3\sqrt{3}-4}{6}\right)\frac{\pi m}{2qB}$。
}

\end{enumerate}

\end{document}
